\section{Background and Related Work}\label{sec:background}

This section presents the concepts of the literature on which Fleet-DT builds and compares Fleet-DT with the related works in four axes: the DTs of vehicles, fleets, and networks (Section~\ref{subsec:bg_dt}), MQTT and the transport of the data of DTs (Section~\ref{subsec:bg_mqtt}), the congestion of a shared wireless link (Section~\ref{subsec:cong}), and the allocation of a shared medium and the regulation of traffic (Section~\ref{subsec:dba}). Section~\ref{subsec:aci} presents the one-step prediction of a measured series, adaptive conformal inference, and the age of information. Table~\ref{tab:related} compares Fleet-DT with the standards and the research systems closest to it.

% Table I of Fleet-DT v2. Requires \usepackage{booktabs,tabularx}.
% Citation keys follow the [key] convention of the revision files; replace by \cite{key} in the IEEE source.
% Every value comes from the cited work (checked 2026-10-15, see REFS_REPORT.md); n.r. marks a value
% that the cited work does not report or that could not be checked. Vehicle DTs without a channel
% allocation (Kremer, Mtowe, Xu, Yan, Zhang, Baniqued, Hartwell) stay in the text of Section II-A/II-B.
\begin{table*}[t]
\caption{Comparison of Fleet-DT with the standards and the research systems closest to it.}
\label{tab:related}
\centering\footnotesize
\setlength{\tabcolsep}{3pt}
\begin{tabularx}{\textwidth}{@{}p{3.0cm}p{2.3cm}p{2.5cm}XXp{2.6cm}@{}}
\toprule
Work & Layer & Allocation unit & Guarantee & Prediction & Validation \\
\midrule
\multicolumn{6}{@{}l}{\emph{Standards}} \\
IEEE 802.11 EDCA [ieee80211std2020] & MAC & Access category & Priority by AIFS, contention window, and TXOP; no bound & None & Standard \\
IEEE 802.11 HCCA [ieee80211std2020] & MAC & Flow (TSPEC) & Admission by the AP; polled transmissions; delay bound declared in the TSPEC & None & Standard \\
IEEE 802.11ax [ieee80211ax2021, khorov2019tutorial] & MAC and PHY & Resource unit per trigger frame; TWT service period & Uplink schedule of the AP & None & Standard \\
IEEE 802.11be [ieee80211be2024, adame2021tsn] & MAC & Restricted TWT service period & Service periods reserved for latency-sensitive traffic & None & Standard \\
IEEE 802.16 [ieee80216std2017] & MAC & Uplink grant per frame (UL-MAP) & Request--grant per frame & None & Standard \\
ETSI DCC [etsi2018dcc] & Access layer & Transmissions of a station & Channel load limited from the CBR measured over 100~ms & None (reacts to the CBR) & Standard \\
\midrule
\multicolumn{6}{@{}l}{\emph{Research systems: medium allocation and admission control}} \\
Kramer et al.\ [kramer2002ipact], 2002 & MAC (EPON) & Window per ONU per polling cycle & Grant capped at a maximum window (limited service) & None (reported queue) & Simulation \\
Rao and Stoica [rao2005overlay], 2005 & Overlay above the IEEE 802.11 MAC & Equal time slot per node & Slots allocated by a distributed algorithm & None & Simulation and Click testbed \\
Yang and Kravets [yang2005contention], 2005 & Network (ad hoc) & Flow & Admission when the nodes in carrier-sensing range have the consumed bandwidth & Available bandwidth from the idle time (moving average) & Analysis and ns-2 simulation \\
Leith et al.\ [leith2005tcp], 2005 & MAC (EDCA parameters) & Traffic class of the AP and of the stations & TCP upload and download fairness & None & Testbed and NS simulation \\
Sarr et al.\ [sarr2008bandwidth], 2008 & Network (ad hoc) & Flow & Admission when the estimate covers the rate of the flow & Available bandwidth from idle overlap and collision probability & n.r. \\
Djukic and Mohapatra [djukic2009soft], 2009 & Software TDMA over commodity IEEE 802.11 & TDMA slot & Clocks synchronized within microseconds; arbitrary TDMA schedule & None & Multi-hop testbed \\
Winstein et al.\ [winstein2013stochastic], 2013 & Transport (end to end) & Bytes per forecast interval & Bounded risk of a delay beyond a target & Stochastic forecast of the link deliveries & Traces of four LTE and 3G networks \\
Wei et al.\ [wei2013rtwifi], 2013 & MAC (TDMA on the IEEE 802.11 TSF) & TDMA slot & Slots at fixed instants & None & Testbed \\
H{\o}iland-J{\o}rgensen et al.\ [hoiland2017ending], 2017 & MAC queues of the AP (Linux) & Airtime per station & Airtime fairness & None & Testbed \\
Shahri et al.\ [shahri2022Extending], 2022 & Application (MQTT) over SDN & MQTT flow & SDN bandwidth reservation; period and deadline per flow & None & Mininet emulation \\
Jaiswal et al.\ [jaiswal2023wiros], 2023 & Application (ROS 2) and EDCA parameters & ROS 2 topic & Desired QoS in a static traffic scenario & None & n.r. \\
\midrule
\multicolumn{6}{@{}l}{\emph{Research systems: network digital twins and digital twin synchronization}} \\
Rusek et al.\ [rusek2020routenet], 2020 & Model of a wired network & Path & None & Delay, jitter, and loss per path (GNN) & Packet-level simulation (OMNeT++) \\
Almasan et al.\ [almasan2022network], 2022 & NDT of a wired network & Routing configuration & None & Delay per path (RouteNet-E) in what-if analysis & Case study against OMNeT++ \\
Almeida et al.\ [almeida2023position], 2023 & NDT of an IEEE 802.11 network (ns-3) & Link & None & Throughput in ns-3 with a learned propagation loss & ns-3 against testbed \\
Li et al.\ [li2024aoi], 2024 & Edge computing & Twin of an object & $1/4$ approximation (static); competitive ratio (dynamic) & None (AoI computed) & Simulation \\
An et al.\ [an2025enabling], 2025 & Application (MQTT) & Update of a UAV twin & None (AoI measured) & None & Simulation \\
\midrule
Fleet-DT & Application (MQTT brokers) above the IEEE 802.11 EDCA & Admission unit (vessel, topic, output) per step & Failure fraction at most $\beta + \alpha_1/(\gamma T) + (1-\beta)(1 + F_{\mathcal{R}}(T))/T$ for every realization of the channel (Theorem~\ref{thm:aci}) & $\hat{\kappa}(t) = (1 - \hat{b}(t))\,\eta(M+1)$ with ACI on the budget & Emulation (ns-3 real time, KVM hosts) \\
\bottomrule
\end{tabularx}
\end{table*}


\subsection{Digital Twins of Vehicles, Fleets, and Networks}\label{subsec:bg_dt}
Kritzinger et al.\ classify virtual representations by the automation of the data flows between a physical object and its representation [kritzinger2018digital]. A digital model exchanges data with the physical object only through manual flows, a digital shadow (DS) receives the data of the physical object automatically, and a DT exchanges data with the physical object automatically in both directions. Grieves and Vickers define the DTP, the DTI, the DTA, and the digital twin environment (DTE) [grieves2017]. The DTP describes a prototypical physical artifact and contains the information necessary to produce a physical version of the artifact. The DTI describes one physical product and remains linked to that product throughout the life of the product. The DTA is the aggregation of the DTIs. The DTE is the integrated multi-domain physics application space in which applications operate on the DTs [grieves2017]. Section~\ref{subsec:dts} defines how Fleet-DT adopts these constructs.

Peng et al.\ [Peng2023] and Wang et al.\ [wang2026Marine] surveyed DTs for unmanned surface and underwater vehicles, and Wang et al.\ classified one of eleven marine robotic cases as a bidirectional DT. Most single-vehicle works implement a DS [kritzinger2018digital]. Bhagavathi et al.\ [Bhagavathi2023] fed sensor data of the Otter autonomous surface vehicle into fault diagnosis models, Berg et al.\ [berg2025Syncing] synchronized the twin of an autonomous surface vehicle by parameter identification with deep reinforcement learning, and Xu et al.\ [xu2025Dynamic] sampled the state of a ship on events to reduce the ship-to-shore traffic. These works couple one twin to one vehicle, leave the transport unspecified, and state no deadline for the update of the twin.

Gervasi et al.\ [gervasi2026Systematic] classified systems of DTs into DTAs, DT networks, federated DTs, composite DTs, and DT ecosystems. Yan et al.\ [yan2026Swarm] and Zhang et al.\ [zhang2024Formation] assigned one twin to each autonomous underwater vehicle of a swarm or to each remotely operated vehicle of a formation and updated the twins periodically. Kim et al.\ [kim2026Command] validated the command-and-control system of five USVs in a DT before the field operation, Wang et al.\ [wang2026Optimizing] trained a cooperative cruising policy for surface vehicles in a DT of the fleet, and Couceiro et al.\ [couceiro2026Paradigm] ran a Unity twin of an aquatic robot swarm with ROS~2 over a Wi-Fi mobile ad hoc network. Outside the marine domain, Baniqued et al.\ [Baniqued2024], Mtowe et al.\ [mtowe2025Edge], and Hartwell et al.\ [Hartwell2024] twinned robot and aero-engine fleets. These twins serve as validation or training environments or as control models. Fleet-DT twins the shared Wi-Fi network of a fleet, with the state of the network of each step in the DTI and the DTIs of the past operations in the DTA.

Zhou et al.\ describe the network digital twin (NDT) as a digital representation whose physical counterpart is a data network, with four elements [zhou2026ndt]. The data hold the historical and real-time state of the network, the models emulate changes in the configuration, the state, or the use of the network elements, the mapping links the NDT to the network in both directions in real time, and the interfaces connect the NDT to the network and to the applications [zhou2026ndt]. Almasan et al.\ represent the network as a graph, use graph neural networks (GNNs) as the models of the NDT, and apply the NDT to what-if analysis, optimization, and troubleshooting [almasan2022network]. RouteNet, the GNN of Rusek et al., takes the topology, the routing, and the traffic matrix of a network and predicts the mean delay, the jitter, and the loss of each path, and the authors trained and tested RouteNet on packet-level simulations of wired topologies with 14 to 50 nodes [rusek2020routenet]. These NDTs model wired networks with routing over several hops. Almeida et al.\ built a DT of a Wi-Fi network in the ns-3 network simulator with a propagation loss model that the authors trained on testbed traces, and this DT runs without a link to the network in operation [almeida2023position]. Fleet-DT twins one Wi-Fi network in operation, in which each vessel reaches the GS in one hop, and predicts, in each step, the capacity of the shared channel from the measurements of this network.

Wu et al.\ survey digital twin networks, in which the twins of physical objects exchange data with these objects over a network [wu2021digital]. Kapteyn et al.\ formulate a DT as a probabilistic graphical model that couples the physical state, the observed data, the digital state, and the control actions over time, and update the digital state of an unmanned aerial vehicle in flight by Bayesian inference [kapteyn2021probabilistic]. The National Academies require continual verification, validation, and uncertainty quantification of a DT for the purpose of the DT [nasem2024foundational]. Fleet-DT states the fidelity of the DT of the network with the metrics of Definition~\ref{def:fidelity} and measures each metric against the ground truth of the emulator (Section~\ref{sec:method}).

\subsection{MQTT and the Transport of the Data of Digital Twins}\label{subsec:bg_mqtt}
In the publish/subscribe protocol MQTT, a client publishes messages on topics to a broker, and the broker delivers each message to the clients that subscribe to the topic of the message [oasis2014mqtt]. MQTT requires a network connection that delivers the bytes in order and without loss, and MQTT implementations use TCP for this connection [oasis2014mqtt]. An MQTT bridge connects two brokers over one client connection and forwards the messages of selected topics between the two brokers, and the configuration of the bridge gives the direction and the QoS level of each forwarded topic [eclipse2026mosquittoconf].

MQTT defines three QoS levels for the delivery of a message from a sender to a receiver [oasis2014mqtt]. With QoS~0, the sender sends the PUBLISH packet once, and the receiver sends no acknowledgment (at most once). With QoS~1, the receiver acknowledges the PUBLISH packet with a PUBACK packet, and the sender keeps the message until the PUBACK packet arrives (at least once). With QoS~2, the sender and the receiver exchange the PUBLISH, PUBREC, PUBREL, and PUBCOMP packets (exactly once). When a client reconnects to a broker with a persistent session (CleanSession flag set to 0), the client and the broker each resend the PUBLISH packets of QoS~1 and~2 that the other side did not acknowledge [oasis2014mqtt]. Each QoS level governs one hop, from a client to a broker or from a broker to a client [oasis2014mqtt].

Zhang et al.\ [Zhang2024] compared DDS, MQTT, and Zenoh under ROS~2 over Ethernet, Wi-Fi, and 4G and measured the lowest delay with Zenoh on the wireless links. Freitas et al.\ [freitas2026Distribution] evaluated MQTT, AMQP, Kafka, DDS, and Zenoh in a DT synchronization architecture and reported that DDS collapsed with 108 nodes under stress, while MQTT and AMQP synchronized states of up to 2~MiB on a degraded network. An et al.\ [an2025enabling] proposed a centralized MQTT framework that updates the DTs of a simulated UAV network with lower latency, jitter, and age of information than MQTT, TCP sockets, and CoAP. Lima et al.\ [lima2025Bridge] showed that the number of MQTT bridges, the packet size, and the topic name affect latency and reliability, and Li et al.\ [li2023Remote] carried sensor data of the research vessel Gunnerus over MQTT to a remote control center.

For video, Herrero [herrero:2023] transported media packets inside MQTT messages, and Kremer et al.\ [kremer2023Teleoperated] and Mtowe et al.\ [mtowe2025Edge] replaced the video stream with a twin of the environment or with edge features to save bandwidth; RTP [rfc3550] carries media in its own session with RTCP feedback. For QoS, Shahri et al.\ [shahri2022Extending] let each MQTT flow declare a period and a deadline and reserved bandwidth through SDN, Palmese et al.\ [palmese2022Adaptive] adjusted the QoS level of each MQTT-SN node at the broker from delay and packet error, and the DDS DEADLINE policy [omg2015Dds] signals missed periodic updates. Jaiswal et al.\ [jaiswal2023wiros] scheduled the flows of the ROS~2 topics of one robot and configured the IEEE~802.11 contention parameters of its Wi-Fi interface for the latency-critical topics. Fleet-DT carries the telemetry, the states, and the video of each vessel over the data bridge of the vessel and admits, in each step, the data of every vessel within the capacity that the DT of the network predicts for the step.

\subsection{Congestion of a Shared Wireless Link and Time-Triggered Communication}\label{subsec:cong}
In IEEE~802.11, the access point and its stations contend for one channel with the distributed coordination function (DCF), so the traffic of both directions shares the capacity of the link [bianchi2000performance]. Bianchi defines the saturation throughput of the DCF as the throughput of the channel when the transmission queue of every station always holds a packet, and shows that the saturation throughput depends on the number of stations that contend for the channel, on the size of the frames, and on the parameters of the physical layer [bianchi2000performance]. Heusse et al.\ [heusse2003performance] show that a station that transmits at a lower bit rate reduces the throughput of every station of the cell and name this effect the performance anomaly of IEEE~802.11. TCP delivers one ordered byte stream per connection [rfc9293], so a segment that waits in the queue delays every later byte of the same connection.

Bianchi models the backoff of each of $n$ saturated stations with a Markov chain and obtains the probability $\tau$ that a station transmits in a slot and the probability $p_{\mathrm{c}}$ that a transmitted frame collides from
\begin{equation}
\tau = \frac{2\,(1 - 2p_{\mathrm{c}})}{(1 - 2p_{\mathrm{c}})(\mathrm{CW}_{\min} + 2) + p_{\mathrm{c}} (\mathrm{CW}_{\min} + 1) \big(1 - (2p_{\mathrm{c}})^{m_{\mathrm{b}}}\big)}, \qquad p_{\mathrm{c}} = 1 - (1 - \tau)^{n-1},
\label{eq:bianchi}
\end{equation}
where $m_{\mathrm{b}}$ is the maximum backoff stage [bianchi2000performance]. With $P_{\mathrm{tr}} = 1 - (1 - \tau)^n$, the probability that a slot holds a transmission, and $P_{\mathrm{s}} = n\tau(1-\tau)^{n-1}/P_{\mathrm{tr}}$, the probability that this transmission succeeds, the fraction of time in which the channel carries successful transmissions is
\begin{equation}
\eta(n) = \frac{P_{\mathrm{tr}} P_{\mathrm{s}}\, T_{\mathrm{succ}}}{(1 - P_{\mathrm{tr}})\,T_{\mathrm{slot}} + P_{\mathrm{tr}} P_{\mathrm{s}}\, T_{\mathrm{succ}} + P_{\mathrm{tr}}(1 - P_{\mathrm{s}})\, T_{\mathrm{coll}}},
\label{eq:eff}
\end{equation}
where $T_{\mathrm{slot}}$ is the slot time, $T_{\mathrm{succ}}$ is the duration of a successful transmission (PHY header, MAC header and payload, SIFS, acknowledgment, and DIFS), and $T_{\mathrm{coll}}$ is the duration of a collision (PHY header, MAC header and payload, and DIFS) [bianchi2000performance]. The saturation throughput of Bianchi is~\eqref{eq:eff} with the transmission time of the mean payload in place of $T_{\mathrm{succ}}$ in the numerator. With EDCA, AIFS of the access category takes the place of DIFS and the contention window of the access category takes the place of $\mathrm{CW}_{\min}$ [ieee80211std2020]. The ns-3 Model Library compares the throughput of the ns-3 Wi-Fi model with the model of Bianchi for IEEE~802.11a, b, and g [nsnam2020wifi].

IEEE~802.11 adds to the DCF the hybrid coordination function (HCF) [ieee80211std2020]. With EDCA, each station keeps one queue per access category, and each access category contends for the channel with its own interframe space and contention window. With HCCA, which IEEE~802.11 defines as optional, a station declares the mean data rate, the maximum burst size, and the delay bound of a flow in a TSPEC, the access point admits or rejects the flow (admission control), and the access point polls the stations of the admitted flows [ieee80211std2020]. In vehicular networks, the ETSI DCC adapts the transmissions of each station to the channel busy ratio (CBR), the fraction of time in which the channel is busy with transmissions, which each station measures over intervals of 100~ms [etsi2018dcc]. The available bandwidth of a link in an interval is the capacity of the link that the current traffic leaves unused in the interval [prasad2003bandwidth].



IEEE~802.11 maps the user priority (UP) of a frame to the access category, with UP~6 and~7 in AC\_VO and UP~0 and~3 in AC\_BE, and RFC~8325 maps the differentiated services code point (DSCP) of an IP packet to the UP [ieee80211std2020, rfc8325]. Leith et al.\ show that, under the DCF, the AP of an infrastructure WLAN receives about the same share of the transmission opportunities as one station, although the AP sends the TCP acknowledgments of every upload and the data of every download, and that EDCA parameters that give the acknowledgments at the AP a shorter AIFS and a smaller $\mathrm{CW}_{\min}$ and the download data at the AP a TXOP of several packets restore the fairness between TCP uploads and downloads [leith2005tcp].

In the time-triggered architecture, the nodes of a distributed system act at instants of a global time base that a clock synchronization establishes, so every node knows in advance the instants at which each message is sent [kopetz2003time]. A firm deadline makes a result useless after the deadline [kopetz2011real].

\subsection{Allocation of a Shared Medium and Traffic Regulation}\label{subsec:dba}
In DBA, a controller divides the capacity of a shared medium among the stations in each cycle from the demand that each station reports. In EPONs, the optical line terminal grants each optical network unit a transmission window in each polling cycle from the queue length that the unit reports, and the limited service of IPACT caps each grant at a maximum window [kramer2002ipact]. In IEEE~802.16, the subscriber stations send bandwidth requests, and the base station grants the uplink bandwidth of each frame in the uplink map [ieee80216std2017]. Both schemes repeat a request--grant cycle, in which the stations report their demand, the controller computes the grants, and the stations transmit within their grants.

Rao and Stoica divide the time into equal slots with loosely synchronized clocks above the IEEE~802.11 MAC and allocate the slots among the nodes with a distributed algorithm [rao2005overlay]. Djukic and Mohapatra synchronize commodity IEEE~802.11 nodes within microseconds and run TDMA schedules over several hops [djukic2009soft]. Wei et al.\ build RT-WiFi, a TDMA MAC on the timing synchronization function of IEEE~802.11, for wireless control loops [wei2013rtwifi]. IEEE~802.11ax lets the AP schedule uplink transmissions on resource units with trigger frames and the service periods of the stations with target wake time (TWT) [ieee80211ax2021, khorov2019tutorial], and IEEE~802.11be adds restricted TWT for latency-sensitive traffic [ieee80211be2024, adame2021tsn]. These mechanisms give slots or resource units to nodes, and each node selects the data of its slot. Fleet-DT admits the data of every vessel within one airtime budget per step in the order of need of the GS, on the EDCA of radios without trigger-based access (Section~\ref{sec:platform}), and Section~\ref{sec:method} compares Fleet-DT with an overlay TDMA and with the slotted variant of Definition~\ref{def:tdma}.

Network calculus bounds the traffic of a flow with an arrival curve, the function that gives, for each duration, the largest number of bytes that the flow sends in an interval of this duration, and a token bucket with a rate and a bucket size gives the arrival curve that starts at the bucket size and grows with the rate [cruz1991calculus, leboudec2001network]. A regulator delays the data of a flow until its output respects an arrival curve, a policer discards or marks the data that exceed an arrival curve, and a strict-priority scheduler serves a class only when the classes of higher priority hold no data [leboudec2001network]. Fleet-DT applies the request--grant cycle above the IEEE~802.11 MAC, to the units of data of the vessels in the order of need of the GS, and regulates the release of the admitted units in each vessel (Section~\ref{sec:dt}).

\subsection{One-Step Prediction and Adaptive Conformal Inference}\label{subsec:aci}
Exponential smoothing with the gain $\lambda \in (0, 1]$ predicts the next value of a measured series $y(t)$ with $\hat{y}(t+1) = (1 - \lambda)\,\hat{y}(t) + \lambda\, y(t)$, and, for a random walk observed with additive noise, a gain that depends on the ratio of the variance of the increments to the variance of the noise minimizes the mean squared error of the one-step prediction [harvey1989forecasting]. TCP predicts the round-trip time of a connection with exponential smoothing and the gain $1/8$ [rfc6298].

Conformal prediction turns a point predictor into a prediction set: with a nonconformity score $S_t$ and the empirical quantile $\hat{Q}_t$ of the past scores, the set $\hat{C}_t(\alpha) = \{y : S_t(y) \leq \hat{Q}_t(1 - \alpha)\}$ covers the next observation with probability at least $1 - \alpha$ when the observations are exchangeable [gibbs2021adaptive]. Adaptive conformal inference (ACI) requires no exchangeability of the observations. With $\mathrm{err}_t = 1$ when the observation of step $t$ falls outside $\hat{C}_t(\alpha_t)$ and $\mathrm{err}_t = 0$ otherwise, ACI updates the level with
\begin{equation}
\alpha_{t+1} = \alpha_t + \gamma\,(\beta - \mathrm{err}_t),
\label{eq:aci_bg}
\end{equation}
where $\beta$ is the target miscoverage and $\gamma > 0$ is the step. The conventions $\hat{Q}_t(x) = \infty$ for $x > 1$ and $\hat{Q}_t(x) = -\infty$ for $x < 0$ keep $\alpha_t \in [-\gamma, 1 + \gamma]$, and summing~\eqref{eq:aci_bg} gives $\big|\frac{1}{T}\sum_{t=1}^{T} \mathrm{err}_t - \beta\big| \leq (\max\{\alpha_1, 1 - \alpha_1\} + \gamma)/(\gamma T)$ for every sequence of observations [gibbs2021adaptive].

The age of information (AoI) of a status at a monitor is the time elapsed since the release of the newest status that the monitor has received [kaul2012real].
Kaul et al.\ show for a first-come first-served queue that the average AoI grows when the update rate approaches zero or the service rate, so an intermediate update rate minimizes the average AoI [kaul2012real]. Yates et al.\ survey the AoI of queues, of scheduled sources, and of sources that share a channel [yates2021age]. Li et al.\ place the twins of mobile objects on edge servers, serve each query from the twin with the lowest AoI, and give a $1/4$ approximation for the static placement [li2024aoi]. These works minimize an average AoI over the updates of the sources. Fleet-DT admits the state of every vessel in every step and bounds the AoI of the state at the GS in every step that meets the hypothesis of Lemma~\ref{lem:delivery} (Corollary~\ref{cor:state}).
